\chapter*{Abstract}

Today, social media has transformed the way people communicate and interact with each other, offering a free space for sharing information, news, photos, and videos. However, free expression and lack of censorship, coupled with users' sense of anonymity, have resulted in an increase in negative and harmful behaviors, such as harassment, discrimination, and bullying, for Internet users. Therefore, this project proposes the design and development of an automatic offensive language detection system that can detect such hate speech and protect users.

The main objective is to offer an accurate, adaptable, and functional real-time solution, initially aimed at digital media outlets that face limitations in opening their comment spaces for fear of offensive content. The system can be integrated as an external service or through an API, allowing content to be filtered before public publication, even in live streaming scenarios. The project considers the design of the classification pipeline, dataset preparation and cleaning, model training and evaluation, as well as a projection of economic viability, technical scalability, and market analysis.

As a result, we hope to contribute to the development of inclusive and accessible tools for the automatic moderation of offensive language in Spanish, promoting safer digital environments that are open to constructive public participation.

\noindent\textbf{Keywords:} Offensive language, natural language processing, automatic content moderation, transformers
